The starting hypothesis was that super El Niño winters are warmer and wetter in Finland than ordinary El Niño winters. The data give a qualified and asymmetric answer. On temperature, the three super winters were about one degree milder than the run of neutral winters of their era in Finland, Sweden, Norway and Denmark in December to February, and about 0.7\,\degC{} milder than ordinary El Niño winters, with fewer frost days, more thaw days and a little less snow. That is the direction the hypothesis predicted, but the effect is within the year-to-year noise of Nordic winters (standard deviation of detrended DJF anomalies around 2\,\degC{} in Finland), it rests on three cases, and it vanishes in January to March. On precipitation, only Finland supports the hypothesis, with a surplus of roughly a quarter of the seasonal total in all three super winters; Sweden and Denmark show none and Iceland the opposite. The absence of any relation between the continuous ONI and either metric means the super group is not the strong end of a dose-response curve.

Three features of the result deserve emphasis. First, the ordinary El Niño group is indistinguishable from neutral winters on every metric and in both windows. If the Nordic signal of El Niño is a late-winter cooling through the stratosphere, as the reviews summarise \cite{bronnimann2007,ineson2009,domeisen2019}, it is too weak to see in 20 winters of national station data; the JFM Finnish value has the predicted sign and nothing more. Second, the super winters differ from ordinary El Niño winters mainly in December to February rather than in late winter, which fits the argument that the strongest events act through a different, more direct route than the moderate ones \cite{toniazzo2006,hardiman2019} and cautions against extrapolating from moderate El Niño composites to the strong events. Third, the three super winters were not mild in absolute terms: Finnish DJF means of $-6$ to $-8$\,\degC{} are ordinary winters, and the 1997/98 winter was slightly colder than the 1991 to 2020 mean. Any statement that a super El Niño brings a mild Finnish winter is a statement about a residual after detrending, and the trend is the larger term.

For practical use, the seasonal forecast literature already treats the tropical Pacific as one predictor among several for European winters \cite{scaife2014}; nothing here suggests that a super El Niño alone should change expectations for a Finnish winter by more than about one degree in DJF, with a wide margin, and by nothing detectable after New Year. The 2023/24 event, which peaked at 1.99\,\degC, sits at the threshold; the companion app shows that lowering the threshold to 1.8 adds it and 1972/73 to the super group and roughly halves the temperature difference. The honest summary is a modest, uncertain early-winter mildness and a Finnish precipitation surplus in three winters, both worth re-examining after the next event rather than acting on now.
